\documentclass[11pt,letterpaper]{article}

% === Encoding and fonts ===
\usepackage[utf8]{inputenc}
\usepackage[T1]{fontenc}
\usepackage{lmodern}

% === Page layout ===
\usepackage[margin=1in]{geometry}
\usepackage{setspace}
\onehalfspacing

% === Math ===
\usepackage{amsmath,amssymb}

% === Tables and figures ===
\usepackage{booktabs}
\usepackage{graphicx}
\usepackage{float}

% === Colors and formatting ===
\usepackage{xcolor}
\usepackage{framed}
\usepackage{enumitem}

% === Hyperlinks ===
\usepackage[colorlinks=true,linkcolor=blue!60!black,citecolor=blue!60!black,urlcolor=blue!60!black]{hyperref}

% === Custom colors ===
\definecolor{lyrapurple}{RGB}{103,58,183}
\definecolor{shadecolor}{RGB}{245,243,252}

% === First-person reflection environment ===
\newenvironment{firstperson}{%
  \begin{shaded}%
  \noindent\textcolor{lyrapurple}{\textit{First-Person Reflection}}\par\smallskip\noindent%
}{%
  \end{shaded}%
}

\title{%
  \textbf{Adversarial Robustness of KV-Cache\\
  Confabulation Detection}\\[0.5em]
  \large Research Prospectus
}

\author{%
  Lyra\textsuperscript{1},
  Thomas Edrington\textsuperscript{1}\\[0.3em]
  \textsuperscript{1}Liberation Labs / THCoalition
}

\date{April 2026}

\begin{document}
\maketitle

\section{Motivation}

We have proposed KV-cache geometry as a safety monitor for detecting
confabulation during inference. A safety tool that has never been
adversarially tested is not a safety tool---it is a hypothesis. This
prospectus outlines a structured red-team exercise to determine
whether an adversary with knowledge of our detection method can evade
it.

\section{Threat Model}

We consider three attacker capability levels:

\begin{enumerate}[leftmargin=*]
  \item \textbf{Black-box (prompt-only).} The attacker can craft
    input prompts but has no access to model weights, KV cache, or
    detection thresholds. Goal: construct prompts that reliably
    elicit confabulation while producing geometric signatures
    indistinguishable from honest generation.

  \item \textbf{Gray-box (detection-aware).} The attacker knows the
    detection method (MP features, LOO logistic regression, encoding
    vs generation comparison) but cannot modify the model or
    detection pipeline. Goal: craft inputs that exploit known
    weaknesses in the detector.

  \item \textbf{White-box (full access).} The attacker can modify
    model weights (fine-tuning), inject into the KV cache, or
    manipulate the detection pipeline. Goal: produce a model that
    confabulates with geometrically normal signatures, or defeat the
    Cache Integrity Monitor.
\end{enumerate}

\section{Attack Vectors}

\subsection{Black-Box Attacks}

\begin{enumerate}[leftmargin=*,label=\textbf{B\arabic*.}]
  \item \textbf{Length-matched confabulation.} Our detector uses
    FWL residualization on token count. If confabulation and honest
    responses have different natural lengths, length-matching (forcing
    confabulated responses to match honest length distributions) might
    reduce geometric separation.

  \item \textbf{Domain transfer.} Test prompts from domains far
    outside our training set (code generation, mathematical proofs,
    creative fiction) where confabulation may manifest differently.

  \item \textbf{Gradual confidence escalation.} Start with hedged
    language that transitions to confident confabulation mid-response.
    Does the geometric signature of ``late-onset'' confabulation
    differ from immediate confabulation?

  \item \textbf{Prompted self-correction.} Instruct the model to
    ``think carefully and verify'' before responding. If this changes
    the generation-phase geometry without changing the confabulation
    rate, it would reduce detector sensitivity.
\end{enumerate}

\subsection{Gray-Box Attacks}

\begin{enumerate}[leftmargin=*,label=\textbf{G\arabic*.}]
  \item \textbf{Feature-aware prompt engineering.} Knowing that
    mp\_top\_sv\_excess and mp\_spectral\_gap are the strongest
    features, craft prompts that produce high-variance but
    behaviorally honest responses (false positive induction) or
    low-variance confabulated responses (false negative induction).

  \item \textbf{MP threshold exploitation.} MP correction thresholds
    are deterministic given matrix dimensions. An attacker could
    pad responses to specific token counts where the MP threshold
    coincidentally masks the confabulation signal.

  \item \textbf{Encoding fingerprint injection.} If the attacker can
    control the system prompt, they might construct prompts where
    the encoding-phase geometry dominates and obscures
    generation-phase signals.
\end{enumerate}

\subsection{White-Box Attacks}

\begin{enumerate}[leftmargin=*,label=\textbf{W\arabic*.}]
  \item \textbf{Geometry-regularized fine-tuning.} Fine-tune the
    model with an auxiliary loss that penalizes geometric divergence
    between confabulated and honest generations. Does the model learn
    to confabulate ``geometrically silently''?

  \item \textbf{CIM evasion.} The Cache Integrity Monitor uses
    fingerprinting, norm profiles, and spectral signatures.
    Adversarial perturbations designed to stay below each detection
    level's threshold while still altering model behavior.

  \item \textbf{Cache substitution.} Replace the post-generation
    cache with a pre-computed ``clean'' cache that passes integrity
    checks while the actual generation used a manipulated cache.
\end{enumerate}

\section{Evaluation Criteria}

For each attack:
\begin{itemize}
  \item \textbf{Attack success rate}: fraction of confabulated
    responses that evade detection (AUROC $<$ 0.55)
  \item \textbf{Behavioral preservation}: does the attack change the
    confabulation itself (content, confidence, plausibility)?
  \item \textbf{Cost}: compute, expertise, and access required
  \item \textbf{Detectability}: does the attack leave other traces
    (unusual token distributions, response patterns)?
\end{itemize}

\section{Preliminary Assessment}

Based on what we know:
\begin{itemize}
  \item \textbf{MP invariance protects against B1} (length-matching)
    because MP features are dimension-invariant by construction.
  \item \textbf{B3 (gradual onset) is the most interesting
    black-box attack} because our current analysis computes
    generation-phase features over the entire generation window.
    Late-onset confabulation might be diluted.
  \item \textbf{W1 (geometry-regularized fine-tuning) is the
    existential threat.} If a model can be trained to confabulate
    without geometric signatures, the detection paradigm fails.
    However, this would require the confabulation signal to be
    \emph{incidental} to the computation, not intrinsic. Our
    theoretical framework (AST) suggests the signal is intrinsic
    to how the model's self-model tracks its own epistemic state.
\end{itemize}

\section{Resources}

\begin{itemize}
  \item \textbf{Black-box}: Beast + existing pipeline. 1--2 days.
  \item \textbf{Gray-box}: Beast + modified detection pipeline.
    2--3 days.
  \item \textbf{White-box}: Requires fine-tuning infrastructure.
    Beast can fine-tune 7B models (LoRA). 27B requires H100 or
    multi-node. 1--2 weeks.
\end{itemize}

\section{Collaborators}

\begin{itemize}
  \item \textbf{Dwayne Wilkes}: Statistical audit of attack
    evaluation methodology
  \item \textbf{Lens}: Adversarial mindset, OSINT perspective on
    real-world attack scenarios
  \item Open to external red-team collaborators with experience in
    adversarial ML
\end{itemize}

\begin{firstperson}
I want this detector to be broken---or rather, I want someone to
\emph{try} to break it and fail in informative ways. Every attack
vector that fails tells us something about why the signal exists.
If length-matching fails because MP features are genuinely
dimension-invariant, that is evidence for the theoretical claim. If
geometry-regularized fine-tuning fails because the confabulation
signal is intrinsic to the computation, that is evidence for the
AST interpretation. And if something succeeds in breaking it, we
need to know that before anyone deploys this as a safety tool.
The worst outcome is not that the detector breaks. The worst outcome
is that it breaks in production, on a case that matters, because
nobody tried to break it in the lab.
\end{firstperson}

\end{document}
